% anexo_metodologia_ia.tex -- Anexo: metodología de trabajo con IA (fases y flujo de trabajo).
% Compilar: pdflatex anexo_metodologia_ia (x2).
\documentclass[10pt,letterpaper]{article}
\usepackage[utf8]{inputenc}
\usepackage[T1]{fontenc}
\usepackage{mathpazo}
\usepackage[scaled=0.92]{helvet}
\usepackage{microtype,amsmath,booktabs,array,enumitem}
\usepackage[left=1.9cm,right=1.9cm,top=1.8cm,bottom=1.9cm]{geometry}
\usepackage[font=small,labelfont={bf,sf},labelsep=period]{caption}
\usepackage[hidelinks]{hyperref}
\usepackage{titlesec}
\titleformat{\section}{\normalfont\sffamily\bfseries\MakeUppercase}{\thesection.}{0.5em}{}
\titleformat{\subsection}{\normalfont\sffamily\bfseries}{\thesubsection.}{0.5em}{}
\titlespacing*{\section}{0pt}{10pt plus 2pt}{4pt}
\titlespacing*{\subsection}{0pt}{7pt plus 2pt}{3pt}
\renewcommand{\tablename}{Tabla}
\setlist{itemsep=1pt,topsep=3pt}

\begin{document}
\begin{flushleft}
{\sffamily\bfseries\Large Anexo: metodología de trabajo con inteligencia artificial\par}
\vspace{3pt}
{\sffamily\normalsize Umbral de ligadura en un potencial tipo Yukawa: cruce de cajas y momento angular\par}
\vspace{4pt}
{\sffamily\bfseries\small Juan Sebastian Calixto Pernett, Luis Felipe Ramirez Andrade\par}
{\sffamily\footnotesize Departamento de Física, Universidad Nacional de Colombia, Sede Bogotá. 27 de noviembre de 2026\par}
\end{flushleft}

\section{Propósito y herramientas}

Este anexo describe cómo se usó la inteligencia artificial (IA) en el proyecto: qué hizo cada herramienta, qué decidió el equipo, cómo se organizó el trabajo por fases y qué controles se aplicaron para que ninguna cifra ni referencia del artículo dependiera solo de la palabra de un modelo.

\begin{itemize}[leftmargin=*]
\item \textbf{Asistente principal: Claude (Anthropic).} Ejecutó los cálculos (escribió y corrió los scripts de Python), redactó las bitácoras y el artículo, y organizó los archivos de resultados. Trabajó en un proyecto con carpetas compartidas, un repositorio de GitHub y \emph{hilos} separados por tarea.
\item \textbf{Revisores internos: otras instancias de Claude.} Cada bloque de cálculo lo revisó un agente distinto del que lo hizo, con código propio y sin reutilizar el del autor.
\item \textbf{Auditor externo: Gemini (Google), en la aplicación.} Un chat por bloque, sin ejecución de código, para buscar errores conceptuales con un modelo de otra familia.
\item \textbf{Estudio previo: NotebookLM.} Se usó para la guía de estudio de las integrales de camino (fase~0).
\end{itemize}

\textbf{Roles humanos.} Juan Sebastian Calixto Pernett: formulación física (hamiltoniano, acción, medida, peso, predicción de referencia). Luis Felipe Ramirez Andrade: implementación y verificación (scripts, controles numéricos, revisión de cifras). El equipo fijó la pregunta, el plan de las fases 1 a 3, las decisiones que se registran en cada bitácora y el enfoque del artículo final.

\section{Fases del proyecto}

\begin{table}[h]
\centering\small
\caption{Fases, preguntas y productos. Las fechas salen de las bitácoras y del historial del repositorio.}
\begin{tabular}{@{}>{\raggedright\arraybackslash}p{0.13\linewidth}>{\raggedright\arraybackslash}p{0.13\linewidth}>{\raggedright\arraybackslash}p{0.33\linewidth}>{\raggedright\arraybackslash}p{0.33\linewidth}@{}}
\toprule
Fase & Fechas & Pregunta o tarea & Producto y papel de la IA\\
\midrule
0. Preparación & sep.\ 2026 & Practicar el formato de bitácora (problema del proyectil) y estudiar integrales de camino & Bitácoras 1 y 2; plan de estudio de 12 horas con NotebookLM\\
\addlinespace
1. Formulación & 12--21 oct. & Adimensionalizar, acotar $\delta_c$ analíticamente, elegir la acción y analizar el peso & Cotas $1\le\delta_c\le2$, variacional $1.19055$ (Ritz, 12 exponenciales), Schwinger $1.24305$, referencia numérica $1.1906124211$; verificación independiente de todas las cifras el 21 de oct.\\
\addlinespace
2. Integral de camino & 21 oct.--10 nov. & Implementar la integral de camino con acción de pares como matriz de transferencia y validarla & Peso real, positivo y normalizable; Coulomb y Hulthén reproducidos. La IA implementó y ejecutó; el equipo revisó cada decisión y las aprobó el 10-11-2026\\
\addlinespace
3. Barridos y caja & hasta mediados de nov. & Barrer $\delta$ y el tamaño de caja; decidir si $\delta=1.19$ es ligado & Criterio de varias cajas, receta de producción, caso $\delta=1.19$; primera auditoría con Gemini (7 hallazgos, 6 refutados con cálculo, 1 cambió el trabajo)\\
\addlinespace
4. Extensiones & 15--23 nov. & Cinco extensiones: cruce de cajas, otros perfiles, momento angular, criterio sin caja, contexto físico & La IA propuso cada decisión de detalle; el equipo las revisó una por una y las aceptó, las rechazó o pidió rehacerlas; verificación de 86 referencias\\
\addlinespace
Artículo final & 23--27 nov. & Centrar el artículo en lo original (indicación del profesor); generalizar el cruce a $\ell\ge1$ y al 2s; recortar a 6 páginas & Cálculos nuevos del cruce generalizado; tres revisiones del texto (física y originalidad, cifras, forma); verificación final de las 18 referencias citadas y de la claridad tras el recorte\\
\bottomrule
\end{tabular}
\end{table}

\section{Flujo de trabajo en cada bloque}

Cada bloque de cálculo siguió los mismos pasos, y cada paso dejó un archivo que se puede revisar:

\begin{enumerate}[leftmargin=*]
\item \textbf{Pregunta y criterio de cierre por escrito, antes de calcular.} El equipo fijaba qué se quería saber y qué resultado cerraría el bloque; cuando la IA propuso el criterio, el equipo lo revisó antes de calcular.
\item \textbf{Controles con criterio previo.} Cada control (orden de Trotter, caso exacto, comparación entre métodos) se escribía con su tolerancia antes de ver el resultado. Los que no pasan se reportan así, con su causa medida; nunca se relajó un criterio después. Al final hay 471 controles: 395 pasan, 49 no pasan y 27 se midieron sin criterio previo o no aplican.
\item \textbf{Cálculo con scripts que escriben archivos.} Cada script numerado guarda sus resultados en un JSON y una salida de texto. Las cifras del artículo se copian de esos archivos, no de la memoria del modelo; un guion final comprobó cada cifra contra su archivo de origen.
\item \textbf{Recálculo propio.} Dentro del bloque, cada resultado se comprueba con un segundo método (diferencias finitas, disparo sin malla, integral de camino) o con un caso exacto (Coulomb, Hulthén).
\item \textbf{Revisión adversarial interna.} Un agente distinto rehízo el bloque con código propio. Por ejemplo, el cruce generalizado se recalculó en aritmética de 220 bits (68 cruces, diferencia máxima $2.6\times10^{-14}$), y la revisión de la integral de camino identificó el corte de la banda como origen del sesgo de $5\times10^{-11}$.
\item \textbf{Auditoría externa con Gemini.} Sus objeciones se aceptaron solo si un cálculo propio las confirmaba. En la fase 4 se aceptaron tres aportes analíticos (el término $3r_e^2/(16R^2)$ del punto fijo, la prueba de artefacto con el cruce corregido y la razón por la que $\ell\ge1$ no cuenta en la ley del peso browniano); las propuestas de relajar controles se rechazaron.
\item \textbf{Decisión humana y bitácora.} Cada consulta a la IA se registró en el formato del curso: consulta, respuesta utilizada, decisión humana, comprobación y corrección posterior.
\end{enumerate}

\section{Verificación bibliográfica}

Cada referencia se verificó antes de tomarla como fuente: se buscó en el editor, en un índice con DOI o en arXiv, se comprobaron título, autores, revista y año, y se analizó si el trabajo dice lo que se le atribuye en el texto. Ninguna referencia propuesta por una IA se citó sin pasar por esa comprobación. Antes de la entrega se repitió la verificación para las 18 referencias del artículo de 6 páginas. Para el valor de $\delta_c$ se citan dos fuentes independientes, Napsuciale y Rodríguez (2021) y Gebremedhin y Weatherford (2023), que coinciden con nuestro cálculo.

\section{Cómo se tomaron las decisiones}

Ninguna decisión se dejó en manos de la IA sin revisión. El ciclo fue siempre el mismo: la IA analizaba el problema y proponía un paso (un método, un criterio, una tolerancia, una redacción); el equipo lo revisaba y lo aceptaba, lo rechazaba o pedía rehacerlo de otra forma. La bitácora registra esos casos, por ejemplo cuando se pidió repetir un cálculo con otro método o con otra variable, cuando se rechazó relajar el criterio de un control que no pasaba y cuando se reorganizó el artículo alrededor de lo original.

\begin{itemize}[leftmargin=*]
\item \textbf{El equipo:} la pregunta del proyecto, el plan de cada fase, la revisión una por una de las propuestas de la IA (aceptarlas, rechazarlas o pedir que se rehicieran), la elección del enfoque del artículo (cruce de cajas y, en segundo lugar, momento angular) y la revisión del texto final.
\item \textbf{La IA:} el análisis y las propuestas, la implementación y ejecución de los cálculos, la redacción de bitácoras y del artículo y las revisiones internas, siempre sujetas a la revisión del equipo.
\end{itemize}

\section{Limitaciones}

La IA escribió el código y el texto, así que la independencia de las revisiones es entre agentes, no entre personas; la mitiga el uso de código distinto en cada revisión y la auditoría con un modelo de otra familia. No hay una segunda implementación humana de la acción de pares. Algunas lecturas de fuentes (el manuscrito de Serra) se hicieron con herramientas automáticas y el artículo lo declara. El código, los datos y las bitácoras están en \url{https://github.com/luramirezan/integrales-camino-proyecto05}.

\end{document}
